\documentclass[aspectratio=169,10pt]{beamer}

% --- math + notation (match the paper) ---
\usepackage{amsmath,amssymb,amsfonts}
\usepackage{bbold}      % \mathbb{1} identity, as in the paper

\usetheme{metropolis}
\metroset{block=fill,numbering=fraction,progressbar=foot}

% ------------------------------------------------------------------
%  Custom 4-stage roadmap header:  Classical -> Body-Fixed -> Quantize -> Quantum H
% ------------------------------------------------------------------
\newcommand{\stagenum}{0}
\newcommand{\setstage}[1]{\renewcommand{\stagenum}{#1}}
\newcommand{\stagelabel}[2]{\ifnum\stagenum=#1 \textbf{\alert{#2}}\else \textcolor{gray}{#2}\fi}
\setbeamertemplate{headline}{%
  \ifnum\stagenum>0%
    \vspace{1.1ex}%
    \hbox to\paperwidth{\hfil
      {\scriptsize \stagelabel{1}{Classical} $\rightarrow$ \stagelabel{2}{Body-Fixed}
       $\rightarrow$ \stagelabel{3}{Quantize} $\rightarrow$ \stagelabel{4}{Quantum $\hat H$}}%
    \hfil}%
  \fi
}

\AtBeginSection[]{%
  \begin{frame}[plain]
    \sectionpage
  \end{frame}}

% shorthand
\newcommand{\Lelec}{\hat{\mathbf{L}}_{\mathrm{elec}}}

\title{Full Molecular Hamiltonian}
\subtitle{From Classical to Quantum: the exact non-relativistic quantization}
\author{Chih-En Shen}
\date{}

\begin{document}

{\setbeamertemplate{headline}{}
\begin{frame}[plain]\titlepage\end{frame}}

% ==================================================================
\begin{frame}{Roadmap: four stages}
\setstage{0}
We build the \alert{exact} non-relativistic molecular Hamiltonian in four stages:
\begin{enumerate}
  \item \textbf{Classical} — lab-frame $H$; separate the centre of mass.
  \item \textbf{Body-Fixed} — co-rotating frame + Lagrangian $\Rightarrow$ classical molecular $H$.
  \item \textbf{Quantize} — Podolsky rule, anomalous algebra, Watson pseudopotential.
  \item \textbf{Quantum $\hat H$} — the exact non-relativistic molecular Hamiltonian.
\end{enumerate}
\vspace{1ex}
\end{frame}

\begin{frame}{Outline}
\setstage{0}
\tableofcontents
\end{frame}

% ==================================================================
\section{Motivation}

\begin{frame}{Why is the molecular Hamiltonian non-trivial?}
\setstage{0}
The lab Hamiltonian couples \emph{all} particles through the Coulomb potential:
\[
H=\sum_j\frac{\mathbf P_j^2}{2M_j}+\sum_i\frac{\mathbf p_i^2}{2m_e}+V .
\]
To do molecular physics we want to separate
\[
\underbrace{\text{translation}}_{\text{COM}}\;\big|\;\underbrace{\text{rotation}}_{\text{Euler angles }\Omega}\;\big|\;\underbrace{\text{vibration}}_{\text{normal coords }Q}\;\big|\;\underbrace{\text{electrons}}_{\bar{\mathbf r}_i}.
\]
\pause
\begin{alertblock}{The catch}
Rotation + vibration are \emph{curvilinear} coordinates. Naively quantizing
$p\to-i\hbar\,\partial$ is \textbf{not} Hermitian there. Getting it right forces
the Podolsky construction, an \emph{anomalous} angular-momentum algebra, and the
Watson pseudopotential.
\end{alertblock}
\end{frame}

\begin{frame}{What we will \emph{not} approximate}
\setstage{0}
We keep everything exact until the very end:
\begin{itemize}
  \item \textbf{No} Born--Oppenheimer / clamped nuclei.
  \item \textbf{No} neglect of \emph{mass polarization}.
  \item \textbf{No} harmonic assumption.
\end{itemize}
\pause
The \emph{only} ingredients are:
\begin{enumerate}
  \item non-relativistic dynamics;
  \item the Eckart frame (a \emph{convention}, not an approximation);
  \item parametrizing the nuclear shape by $3N_{\mathrm{nucl}}-6$ normal coordinates (exact).
\end{enumerate}
\vspace{0.5ex}
\centering\itshape The endpoint is the exact quantum $\hat H_{\mathrm{rve}}$ — the rigorous starting point \\ for \emph{any} later approximation.
\end{frame}

% ==================================================================
\setstage{1}
\section{Classical Cartesian Hamiltonian}

\begin{frame}{The starting point (lab frame)}
\setstage{1}
$N_{\mathrm{nucl}}$ nuclei ($M_j,Z_je,\mathbf R_j,\mathbf P_j$) and $N_{\mathrm{elec}}$ electrons ($m_e,-e,\mathbf r_i,\mathbf p_i$), all in the space-fixed frame:
\begin{equation*}
H = \underbrace{\sum_{j}\frac{\mathbf{P}_j^{\,2}}{2M_j}}_{T_N} + \underbrace{\sum_{i}\frac{\mathbf{p}_i^{\,2}}{2m_e}}_{T_e} + V, \tag{1.1}
\end{equation*}
\begin{equation*}
V = \frac{1}{4\pi\epsilon_0}\left(\frac12\sum_{i\ne j}\frac{Z_i Z_j e^2}{R_{ij}} - \sum_{i,j}\frac{Z_j e^2}{|\mathbf{R}_j-\mathbf{r}_i|} + \frac12\sum_{i\ne j}\frac{e^2}{r_{ij}}\right). \tag{1.2}
\end{equation*}
Canonical brackets $\{R_{j\alpha},P_{k\beta}\}=\delta_{jk}\delta_{\alpha\beta}$, $\{r_{i\alpha},p_{j\beta}\}=\delta_{ij}\delta_{\alpha\beta}$.
\vspace{0.5ex}

\centering\footnotesize Everything that follows is an exact rewriting of (1.1)--(1.2).
\end{frame}

% ==================================================================
\setstage{1}
\section{Nuclear centre-of-mass separation}

\begin{frame}{New coordinates: nuclear COM + relative}
\setstage{1}
Split off the nuclear centre of mass and refer electrons to it:
\begin{align*}
\mathbf{R}_{\mathrm{cm}} &= \frac{1}{M_N}\sum_j M_j\mathbf{R}_j,\qquad M_N=\sum_j M_j, \tag{2.1}\\
\mathbf{R}'_j &= \mathbf{R}_j-\mathbf{R}_{\mathrm{cm}},\qquad \mathbf{r}'_i=\mathbf{r}_i-\mathbf{R}_{\mathrm{cm}}. \tag{2.2}
\end{align*}
The relative nuclear vectors are linearly dependent:
\begin{equation*}
\sum_j M_j\mathbf{R}'_j=0. \tag{2.3}
\end{equation*}
\pause
A point ($F_2$-)transformation gives the conjugate momenta [dual constraint $\sum_k\mathbf P'_k=0$]:
\begin{equation*}
\boxed{\;\mathbf{P}_j = \tfrac{M_j}{M_N}\bigl(\mathbf{P}_{\mathrm{cm}}-\textstyle\sum_i\mathbf{p}_i\bigr)+\mathbf{P}'_j,\qquad \mathbf{p}_i=\mathbf{p}'_i\;}\tag{2.7}
\end{equation*}
with $\mathbf P_{\mathrm{cm}}=\sum_j\mathbf P_j+\sum_i\mathbf p_i$ the \emph{total} momentum.
\end{frame}

\begin{frame}{Hamiltonian after the transformation}
\setstage{1}
Insert (2.7) and expand the nuclear kinetic energy, keeping $\mathbf P_{\mathrm{cm}}$ \emph{arbitrary}:
\begin{equation*}
H = \underbrace{\frac{\mathbf{P}_{\mathrm{cm}}^{\,2}}{2M_N}}_{(A)\ \text{wrong mass}}\underbrace{-\frac{\mathbf{P}_{\mathrm{cm}}\!\cdot\!\sum_i\mathbf{p}_i}{M_N}}_{(B)\ \text{cross-coupling}}+\underbrace{\frac{(\sum_i\mathbf{p}_i)^2}{2M_N}}_{\text{mass polarization}}+\sum_j\frac{\mathbf{P}_j^{'2}}{2M_j}+\sum_i\frac{\mathbf{p}_i^{\,2}}{2m_e}+V. \tag{2.10}
\end{equation*}
\begin{block}{Why the mass-polarization term?}
It is the \emph{price} of referring electrons to the \emph{nuclear} COM (not the total COM): the electron momenta no longer combine additively with $M_N$. We keep it exactly --- it is never dropped.
\end{block}
\pause
Artifacts (A) and (B) appear only because the split used the \emph{nuclear} COM. We remove both \alert{exactly, for any $\mathbf P_{\mathrm{cm}}$} --- with \emph{no} need to set $\mathbf P_{\mathrm{cm}}=0$.
\end{frame}

\begin{frame}{Eliminating the cross-coupling, for any $\mathbf P_{\mathrm{cm}}$}
\setstage{1}
A single classical canonical transformation generated by
$G=\tfrac{m_e}{M_{\mathrm{tot}}}\mathbf{P}_{\mathrm{cm}}\!\cdot\!\sum_i\mathbf{r}_i'$ removes \emph{both} artifacts at once:
\begin{equation*}
\tilde{\mathbf r}'_i=\mathbf r'_i,\quad \tilde{\mathbf p}'_i=\mathbf p'_i-\tfrac{m_e}{M_{\mathrm{tot}}}\mathbf P_{\mathrm{cm}},\quad
\tilde{\mathbf R}_{\mathrm{cm}}=\mathbf R_{\mathrm{cm}}+\tfrac{m_e}{M_{\mathrm{tot}}}\textstyle\sum_i\mathbf r'_i. \tag{2.15}
\end{equation*}
\begin{block}{Why this works}
It relabels the COM from \emph{nuclear} to \emph{total} --- a Galilean boost. The wrong mass (A) is corrected $M_N\!\to\!M_{\mathrm{tot}}$ and the cross-coupling (B) is killed, \emph{exactly and for any} $\mathbf P_{\mathrm{cm}}$. So we never assume $\mathbf P_{\mathrm{cm}}=0$.
\end{block}
\end{frame}

\begin{frame}{The molecular Hamiltonian $\tilde H$: correct mass, no cross-coupling}
\setstage{1}
Exact for \emph{any} $\mathbf P_{\mathrm{cm}}$, $\tilde H$ splits into a free COM term and the \alert{rotation--vibration--electronic} Hamiltonian $H_{\mathrm{rve}}$:
\begin{equation*}
\boxed{\;\tilde H = \frac{\mathbf{P}_{\mathrm{cm}}^{\,2}}{2M_{\mathrm{tot}}}+H_{\mathrm{rve}},\qquad
H_{\mathrm{rve}}=\sum_j\frac{\mathbf{P}_j^{'2}}{2M_j}+\sum_i\frac{\tilde{\mathbf{p}}_i^{'2}}{2m_e}+\frac{1}{2M_N}\Bigl(\sum_i\tilde{\mathbf{p}}'_i\Bigr)^{\!2}+V.\;} \tag{2.22}
\end{equation*}
\begin{itemize}
  \item COM term carries the \emph{correct} total mass $M_{\mathrm{tot}}=M_N+N_{\mathrm{elec}}m_e$; $\mathbf P_{\mathrm{cm}}$ is a decoupled, conserved spectator.
  \item $H_{\mathrm{rve}}$ = nuclear KE $+$ electron KE $+$ mass polarization $+$ Coulomb.
\end{itemize}
\vspace{0.5ex}
\centering\itshape $H_{\mathrm{rve}}$ is what we carry into the body-fixed frame.
\end{frame}

% ==================================================================
\setstage{2}
\section{Body-fixed frame and the Lagrangian}

\begin{frame}{Body-fixed frame: rotate with the nuclei}
\setstage{2}
Introduce a rotation $\mathbf S(\Omega)$ (Euler angles $\Omega=\phi,\theta,\chi$) and body-fixed (BF) coordinates:
\begin{equation*}
\mathbf{R}'_j=\mathbf{S}(\Omega)\,\bar{\mathbf{R}}_j(Q),\qquad \mathbf{r}'_i=\mathbf{S}(\Omega)\,\bar{\mathbf{r}}_i. \tag{3.1}
\end{equation*}
The \emph{Eckart} conditions fix $\mathbf S$ from the \emph{nuclei}; the nuclear shape is rectilinear mass-weighted \alert{normal coordinates} $Q_b$:
\begin{equation*}
\bar{\mathbf{R}}_j(Q)=\bar{\mathbf{R}}_j^{\mathrm{eq}}+\frac{1}{\sqrt{M_j}}\sum_b\mathbf{l}_{jb}\,Q_b ,
\qquad \sum_j\mathbf{l}_{ja}\!\cdot\!\mathbf{l}_{jb}=\delta_{ab}. \tag{3.1a}
\end{equation*}
The BF angular velocity $\boldsymbol{\omega}$ is defined by
\begin{equation*}
(\mathbf{S}^{T}\dot{\mathbf{S}})_{\alpha\beta}=-\epsilon_{\alpha\beta\gamma}\,\omega_\gamma . \tag{3.2}
\end{equation*}
\end{frame}

\begin{frame}{A Lagrangian that keeps electrons in momentum form}
\setstage{2}
Start from $\tilde H$ (2.22). We need an action principle for the body-fixed \emph{nuclear} coordinates --- but the electron kinetic energy + mass polarization
$H_e=\sum_i\tilde{\mathbf p}_i'^{\,2}/2m_e+\tfrac{1}{2M_N}(\sum_i\tilde{\mathbf p}'_i)^2$
is a \emph{rotational scalar}. So we Legendre-transform only the COM and nuclear momenta and \alert{keep the electrons in momentum form}, giving the Lagrangian
\begin{equation*}
\boxed{\;\mathcal{L}=\tfrac12 M_{\mathrm{tot}}|\dot{\tilde{\mathbf{R}}}_{\mathrm{cm}}|^2+\tfrac12\sum_j M_j|\dot{\mathbf{R}}_j'|^2-H_e(\tilde{\mathbf{p}}'_i)-V.\;}\tag{3.6}
\end{equation*}
\begin{block}{Why keep electrons in momentum form?}
It avoids converting the electron kinetic energy to velocity form and back: the electrons ride along as the scalar $H_e$. So $\mathcal L$ is a Lagrangian in the nuclear/COM coordinates $(\tilde{\mathbf R}_{\mathrm{cm}},\mathbf R'_j)$ and carries the electrons through their momenta $(\mathbf r'_i,\tilde{\mathbf p}'_i)$.
\end{block}
\end{frame}

\begin{frame}{Body-fixed Lagrangian}
\setstage{2}
$H_e$ is built from $\tilde{\mathbf p}'_i\!\cdot\!\tilde{\mathbf p}'_l$; under the passive rotation $\bar{\mathbf p}_i\equiv\mathbf S^{T}\tilde{\mathbf p}'_i$ it is \emph{unchanged} in form. Substituting the BF velocities $\dot{\mathbf R}'_j=\mathbf S(\boldsymbol{\omega}\!\times\!\bar{\mathbf R}_j+\dot{\bar{\mathbf R}}_j)$,
\begin{equation*}
\mathcal{L}=\tfrac12 M_{\mathrm{tot}}|\dot{\tilde{\mathbf{R}}}_{\mathrm{cm}}|^2
+\tfrac12\!\sum_j M_j|\boldsymbol{\omega}\!\times\!\bar{\mathbf{R}}_j+\dot{\bar{\mathbf{R}}}_j|^2
-\sum_i\frac{\bar{\mathbf{p}}_i^{\,2}}{2m_e}-\frac{1}{2M_N}\Bigl(\sum_i\bar{\mathbf{p}}_i\Bigr)^{2}-V. \tag{3.7}
\end{equation*}
Only the \emph{nuclear} kinetic term carries $\boldsymbol{\omega}$; $V$ is now manifestly $\mathbf S$-independent.
\end{frame}

\begin{frame}{Expand the nuclear term}
\setstage{2}
The vector identity (for any $\mathbf a,\dot{\mathbf a},\boldsymbol\omega$)
\begin{equation*}
\tfrac{m}{2}|\boldsymbol{\omega}\!\times\!\mathbf{a}+\dot{\mathbf{a}}|^{2}
=\underbrace{\tfrac{m}{2}\boldsymbol{\omega}^{T}(a^{2}\mathbb{1}-\mathbf{a}\mathbf{a}^{T})\boldsymbol{\omega}}_{\text{rotational}}
+\underbrace{m\,\boldsymbol{\omega}\!\cdot\!(\mathbf{a}\!\times\!\dot{\mathbf{a}})}_{\text{Coriolis}}
+\underbrace{\tfrac{m}{2}|\dot{\mathbf{a}}|^{2}}_{\text{vibrational}} \tag{3.9c}
\end{equation*}
\pause
applied to the nuclei ($m\!\to\!M_j$, $\mathbf a\!\to\!\bar{\mathbf R}_j$) gives the assembled Lagrangian
\begin{equation*}
\boxed{\;\mathcal{L}=\tfrac12 M_{\mathrm{tot}}|\dot{\tilde{\mathbf{R}}}_{\mathrm{cm}}|^2
+\underbrace{\tfrac12\boldsymbol{\omega}^{T}\mathbf{I}_N\boldsymbol{\omega}+\boldsymbol{\omega}\!\cdot\!\mathbf{L}_N^{\mathrm{vib}}+T_{\mathrm{vib}}^{N}}_{\text{nuclear}}-H_e-V.\;} \tag{3.13}
\end{equation*}
\alert{The rotational kernel is the \emph{nuclear} inertia $\mathbf I_N$ alone} --- no electronic inertia, because the electrons were never expanded in $\boldsymbol\omega$.
\end{frame}

\begin{frame}{Canonical momenta and the total angular momentum}
\setstage{2}
\textbf{Nuclear} angular momentum (from $\partial\mathcal L/\partial\boldsymbol\omega$):
\begin{equation*}
\mathbf{J}_N=\mathbf{I}_N\boldsymbol{\omega}+\mathbf{L}_N^{\mathrm{vib}}. \tag{3.15}
\end{equation*}
The BF frame rotates nuclei \emph{and} electrons, so the \alert{total} angular momentum is the sum [$\mathbf L_{\mathrm{elec}}=\sum_i\bar{\mathbf r}_i\!\times\!\bar{\mathbf p}_i$]:
\begin{equation*}
\mathbf{J}=\mathbf{J}_N+\mathbf{L}_{\mathrm{elec}}=\mathbf{I}_N\boldsymbol{\omega}+\mathbf{L}_N^{\mathrm{vib}}+\mathbf{L}_{\mathrm{elec}}. \tag{3.16}
\end{equation*}
Vibrational momentum, with the \alert{Coriolis vectors} $\boldsymbol\sigma_b$:
\begin{equation*}
P_b=\boldsymbol{\omega}\!\cdot\!\boldsymbol{\sigma}_b+\dot Q_b,\qquad
\boldsymbol{\sigma}_b=\sum_a\boldsymbol{\zeta}_{ab}Q_a,\quad \boldsymbol{\zeta}_{ab}=\sum_j\mathbf{l}_{ja}\!\times\!\mathbf{l}_{jb}. \tag{3.17a,b}
\end{equation*}
\end{frame}

\begin{frame}{The key identity}
\setstage{2}
Subtract the electronic part from (3.16):
\begin{equation*}
\boxed{\;\mathbf{J}-\mathbf{L}_{\mathrm{elec}}=\mathbf{I}_N(Q)\,\boldsymbol{\omega}+\mathbf{L}_N^{\mathrm{vib}}.\;} \tag{3.19}
\end{equation*}
\begin{block}{Why this matters}
The right side is \emph{purely nuclear}: the rotational kernel is $\mathbf I_N$, \textbf{not} the full inertia. All electronic + mass-polarization content sits on the left, in the kinematic shift $\mathbf J\to\mathbf J-\mathbf L_{\mathrm{elec}}$.
\end{block}
\pause
\footnotesize\emph{Cross-check:} the velocity-form route instead generates the electronic inertia $\mathbf I_e,\mathbf I_X$ in \emph{both} $\mathbf J$ and $\mathbf L_{\mathrm{elec}}$ --- they cancel exactly, reproducing (3.19). The Lagrangian makes the cancellation automatic.
\end{frame}

\begin{frame}{Watson reciprocal inertia}
\setstage{2}
Invert (3.19) for $\boldsymbol\omega$: write $\mathbf L_N^{\mathrm{vib}}=\boldsymbol\pi-\boldsymbol\sigma\boldsymbol\sigma^{\mathsf T}\boldsymbol\omega$ with
\begin{equation*}
\boldsymbol{\pi}=\sum_b\boldsymbol{\sigma}_bP_b\ \ (\text{field-free vibrational ang.\ mom.}),\qquad
\boldsymbol{\sigma}\boldsymbol{\sigma}^{\mathsf T}=\sum_b\boldsymbol{\sigma}_b\boldsymbol{\sigma}_b^{\mathsf T}. \tag{3.20a}
\end{equation*}
Then $\mathbf J-\mathbf L_{\mathrm{elec}}=(\mathbf I_N-\boldsymbol\sigma\boldsymbol\sigma^{\mathsf T})\boldsymbol\omega+\boldsymbol\pi$, so
\begin{equation*}
\boxed{\;\boldsymbol{\omega}=\boldsymbol{\mu}\bigl(\mathbf{J}-\mathbf{L}_{\mathrm{elec}}-\boldsymbol{\pi}\bigr),\qquad
\boldsymbol{\mu}\equiv(\mathbf{I}_N-\boldsymbol{\sigma}\boldsymbol{\sigma}^{\mathsf T})^{-1}.\;}\tag{3.20}
\end{equation*}
$\boldsymbol\mu$ is the \alert{Watson reciprocal-inertia tensor} --- built from \emph{nuclear} quantities only, with no Wilson $G$-matrix.
\end{frame}

\begin{frame}{The exact classical molecular Hamiltonian}
\setstage{2}
Legendre back; the rovibrational kinetic energy becomes $\tfrac12(\mathbf J-\mathbf L_{\mathrm{elec}}-\boldsymbol\pi)^{\mathsf T}\boldsymbol\mu(\mathbf J-\mathbf L_{\mathrm{elec}}-\boldsymbol\pi)+\tfrac12\sum_bP_b^2$:
\begin{equation*}
\boxed{
\begin{aligned}
H_{\mathrm{rve}}=\;&\tfrac12\bigl(\mathbf{J}-\mathbf{L}_{\mathrm{elec}}-\boldsymbol{\pi}\bigr)^{\mathsf T}\boldsymbol{\mu}(Q)\bigl(\mathbf{J}-\mathbf{L}_{\mathrm{elec}}-\boldsymbol{\pi}\bigr)+\tfrac12\!\sum_bP_b^2\\
&+\sum_i\frac{\bar{\mathbf p}_i^{\,2}}{2m_e}+\frac{1}{2M_N}\Bigl(\sum_i\bar{\mathbf p}_i\Bigr)^{2}+V_{NN}+V_{Ne}+V_{ee}.
\end{aligned}}\tag{3.22}
\end{equation*}
\centering Exact, classical, in BF + Eckart coordinates. \alert{Now we quantize.}
\end{frame}

% ==================================================================
\setstage{3}
\section{The quantization problem}

\begin{frame}{Why naive quantization fails in curvilinear coordinates}
\setstage{3}
Label the configuration coordinates and write the classical kinetic energy as a quadratic form:
\begin{equation*}
\{q^A\}=\{\phi,\theta,\chi,\;Q_1,\dots,Q_{3N_{\mathrm{nucl}}-6},\;\bar{\mathbf r}_1,\dots\},\qquad
T=\tfrac12\,g^{AB}(q)\,\pi_A\pi_B. \tag{4.1, 4.2}
\end{equation*}
The metric $g_{AB}$ inherits the block structure of (3.22): Watson $\boldsymbol\mu(Q)$ (rotation), identity (vibration), Cartesian (electrons).
\begin{alertblock}{The problem}
The Euler angles and $Q$ are \emph{curvilinear}. The natural Hilbert space is $L^2\!\bigl(\sqrt{|g|}\,d^nq\bigr)$, but $\hat\pi_A=-i\hbar\partial_A$ makes $\tfrac12 g^{AB}\hat\pi_A\hat\pi_B$ \textbf{non-Hermitian} there ($g^{AB}$ depends on $q$ and does not commute with $\partial_A$).
\end{alertblock}
\end{frame}

% ==================================================================
\section{Podolsky quantization}

\begin{frame}{The Podolsky / Laplace--Beltrami prescription}
\setstage{3}
The unique kinetic operator Hermitian on $L^2(\mathcal C,\sqrt{|g|}\,d^nq)$ is the Laplace--Beltrami operator:
\begin{equation*}
\boxed{\;\hat T=-\frac{\hbar^2}{2}\,\frac{1}{\sqrt{|g|}}\,\partial_A\!\bigl(\sqrt{|g|}\,g^{AB}\,\partial_B\bigr).\;}\tag{4.3}
\end{equation*}
Equivalently, with the \emph{Hermitian} momenta $\hat\pi_A=-i\hbar\bigl(\partial_A+\tfrac12\partial_A\ln\sqrt{|g|}\bigr)$, one has $\hat T=\tfrac12\hat\pi_A g^{AB}\hat\pi_B$.
\begin{block}{Why this form?}
The $\sqrt{|g|}$ factors symmetrize the operator about the curved measure. The price: a residual $Q$-dependent $c$-number --- the \alert{Watson pseudopotential} --- once we flatten the measure.
\end{block}
\end{frame}

\begin{frame}{The molecular volume element}
\setstage{3}
Building $g_{AB}$ from the velocity-form kinetic energy and reducing the determinant (Schur complements over the electron and vibrational blocks) gives
\begin{equation*}
\boxed{\;\sqrt{|g|}\,d^nq=\underbrace{\sin\theta\,d\phi\,d\theta\,d\chi}_{\text{Haar on }SO(3)}\cdot\underbrace{\sqrt{\det\mathbf I'(Q)}\,\textstyle\prod_b dQ_b}_{\text{rot.--vib.\ Jacobian}}\cdot\underbrace{\textstyle\prod_i d^3\bar{\mathbf r}_i}_{\text{flat}}.\;}\tag{4.4}
\end{equation*}
The Jacobian carries the \alert{Watson modified inertia} $\mathbf I'=\mathbf I_N-\boldsymbol\sigma\boldsymbol\sigma^{\mathsf T}$, \emph{not} the bare $\mathbf I_N$: the rotation--vibration Coriolis coupling reduces $\det\mathbf I_N\to\det\mathbf I'$.
\end{frame}

\begin{frame}{Rescaling to a flat measure}
\setstage{3}
To work on the simpler measure $d\mu_W=\sin\theta\,d\phi d\theta d\chi\prod_b dQ_b\prod_i d^3\bar{\mathbf r}_i$, rescale the wavefunction by the $Q$-Jacobian $\gamma\equiv\det\mathbf I'$:
\begin{equation*}
\psi_W\equiv\bigl(\det\mathbf I'(Q)\bigr)^{1/4}\psi,\qquad
\hat T_W=\mathcal J^{1/2}\hat T\,\mathcal J^{-1/2},\quad \mathcal J=(\det\mathbf I')^{1/2}. \tag{4.5}
\end{equation*}
\begin{block}{Why rescale?}
$\psi_W\in L^2(d\mu_W)$ with the \emph{flat} vibrational measure, so $\hat P_b=-i\hbar\partial_{Q_b}$ is Hermitian. The similarity does not commute through $\gamma(Q)$ --- the leftover is exactly the pseudopotential $\hat U(Q)$.
\end{block}
\end{frame}

% ==================================================================
\section{Anomalous body-fixed angular momentum}

\begin{frame}{Body-fixed angular momentum has a sign flip}
\setstage{3}
Total $\mathbf J$ is an SF observable; its BF components are projections through $\mathbf S(\Omega)$:
\begin{equation*}
\hat J_\alpha^{\mathrm{BF}}=\mathbf S_{\beta\alpha}\,\hat J_\beta^{\mathrm{SF}}. \tag{4.6}
\end{equation*}
But $\mathbf S(\Omega)$ depends on the Euler angles that $\hat J^{\mathrm{SF}}$ differentiates. Commuting $\hat J^{\mathrm{SF}}$ past $\mathbf S$ adds terms that \emph{outweigh and flip} the usual algebra:
\begin{equation*}
\boxed{\;[\hat J_\alpha^{\mathrm{BF}},\hat J_\beta^{\mathrm{BF}}]=-\,i\hbar\,\epsilon_{\alpha\beta\gamma}\,\hat J_\gamma^{\mathrm{BF}}.\;}\tag{4.7}
\end{equation*}
\centering The famous \alert{anomalous commutator} (Klein; Van Vleck): note the $\boldsymbol{-}$ sign.
\end{frame}

\begin{frame}{Normal vs.\ anomalous; commuting sectors}
\setstage{3}
The electronic angular momentum is built intrinsically from BF electron operators --- \emph{normal} algebra:
\begin{equation*}
[\hat L_{\mathrm{elec},\alpha},\hat L_{\mathrm{elec},\beta}]=+\,i\hbar\,\epsilon_{\alpha\beta\gamma}\,\hat L_{\mathrm{elec},\gamma},\qquad
[\hat J_\alpha^{\mathrm{BF}},\hat L_{\mathrm{elec},\beta}]=0. \tag{4.8, 4.9}
\end{equation*}
\begin{block}{Why it (barely) matters}
The sign flip controls the ordering of $\hat J_\alpha\hat J_\beta$ in the kinetic operator. But contracted with the \emph{symmetric} $\mu_{\alpha\beta}$ the antisymmetric commutator drops out ($\mu_{\alpha\beta}\epsilon_{\alpha\beta\gamma}=0$) --- so it contributes \emph{no} extra potential. We use this next.
\end{block}
\end{frame}

% ==================================================================
\section{Watson pseudopotential}

\begin{frame}{The general curvilinear extra-potential}
\setstage{3}
Conjugating the Laplace--Beltrami operator $\hat T_{\mathrm{LB}}=-\tfrac{\hbar^2}{2}g^{-1/2}\partial_A(g^{1/2}g^{AB}\partial_B)$ by $g^{1/4}$ and collecting terms (first-order pieces cancel by symmetry of $g^{AB}$) gives the exact, coordinate-general result
\begin{equation*}
\boxed{\;\hat T_W=\tfrac12\,\hat p_A\,g^{AB}\,\hat p_B+U,\qquad
U=\frac{\hbar^2}{8}\Bigl[\tfrac14 g^{AB}(\partial_A\ln g)(\partial_B\ln g)+\partial_A(g^{AB}\partial_B\ln g)\Bigr].\;}\tag{4.10d}
\end{equation*}
\begin{block}{Why $\hbar^2/8$?}
The two half-powers $g^{\pm1/4}$ of the rescaling each contribute $\tfrac12\partial\ln g$; their product/derivative is $\mathcal O(\hbar^2)$ with prefactor $\tfrac18$.
\end{block}
\end{frame}

\begin{frame}{The molecular metric and its determinant}
\setstage{3}
In the body-frame momentum basis $(\hat J_\alpha,\hat P_b)$ [with $\boldsymbol\pi=\boldsymbol\sigma^{\mathsf T}\hat P$], read off
\begin{equation*}
g^{AB}=\begin{pmatrix}\boldsymbol\mu & -\boldsymbol\mu\boldsymbol\sigma^{\mathsf T}\\[2pt] -\boldsymbol\sigma\boldsymbol\mu & \mathbb 1+\boldsymbol\sigma\boldsymbol\mu\boldsymbol\sigma^{\mathsf T}\end{pmatrix},
\qquad \det g^{AB}=\det\boldsymbol\mu \tag{4.10e,f}
\end{equation*}
(Schur complement of the $\boldsymbol\mu$ block). Hence $g=1/\det\boldsymbol\mu=\det\mathbf I'\equiv\gamma(Q)$. Since the $\hat J_\alpha$ already carry $\sin\theta$, only $\gamma(Q)$ is rescaled: set $\ln g\to\ln\gamma$ in (4.10d).
\end{frame}

\begin{frame}{Rotational sector: no $c$-number}
\setstage{3}
The only ordering ambiguity is in $\hat J_\alpha\hat J_\beta$. Split and use the anomalous algebra (4.7):
\begin{equation*}
\tfrac12\mu_{\alpha\beta}\hat J_\alpha\hat J_\beta
=\tfrac14\mu_{\alpha\beta}\{\hat J_\alpha,\hat J_\beta\}
-\frac{i\hbar}{4}\,\mu_{\alpha\beta}\,\epsilon_{\alpha\beta\gamma}\,\hat J_\gamma. \tag{4.10g}
\end{equation*}
Symmetric $\mu_{\alpha\beta}$ $\times$ antisymmetric $\epsilon_{\alpha\beta\gamma}=0$ $\Rightarrow$ the commutator term vanishes. \alert{The rotational kernel is Hermitian as written, with no extra potential.}
\end{frame}

\begin{frame}{Vibrational sector: the rescaling, step by step}
\setstage{3}
With $g^{ab}=\delta_{ab}$ but measure $\gamma(Q)$, the Podolsky vibrational operator is $\hat T_{\mathrm{vib}}=-\tfrac{\hbar^2}{2}\gamma^{-1/2}\sum_b\partial_{Q_b}(\gamma^{1/2}\partial_{Q_b})$. Conjugate by $\gamma^{1/4}$, working outward (one product rule per line):
\begin{align*}
\partial_{Q_b}(\gamma^{-1/4}\psi_W)&=\gamma^{-1/4}\bigl[\partial_{Q_b}\psi_W-\tfrac14(\partial_{Q_b}\!\ln\gamma)\psi_W\bigr],\\
\gamma^{1/2}\partial_{Q_b}(\gamma^{-1/4}\psi_W)&=\gamma^{1/4}\bigl[\partial_{Q_b}\psi_W-\tfrac14(\partial_{Q_b}\!\ln\gamma)\psi_W\bigr],\\
\partial_{Q_b}\bigl(\gamma^{1/4}[\cdots]\bigr)&=\gamma^{1/4}\bigl[\partial_{Q_b}^2\psi_W-\tfrac{1}{16}(\partial_{Q_b}\!\ln\gamma)^2\psi_W-\tfrac14(\partial_{Q_b}^2\ln\gamma)\psi_W\bigr].
\end{align*}
The cross terms $\pm\tfrac14(\partial\ln\gamma)\partial\psi_W$ cancel. Result:
\begin{equation*}
\gamma^{1/4}\hat T_{\mathrm{vib}}\gamma^{-1/4}=\tfrac12\sum_b\hat P_b^2+U_{\mathrm{vib}},\quad
U_{\mathrm{vib}}=\frac{\hbar^2}{8}\sum_b\Bigl[\partial_{Q_b}^2\ln\gamma+\tfrac14(\partial_{Q_b}\!\ln\gamma)^2\Bigr]. \tag{4.10k}
\end{equation*}
\end{frame}

\begin{frame}{Collapse to the Watson trace}
\setstage{3}
Two exact identities finish the job. \textbf{Jacobi:}
\begin{equation*}
\partial_b\ln\gamma=\operatorname{tr}(\boldsymbol\mu\,\partial_b\mathbf I'),\qquad \partial_b\boldsymbol\mu=-\boldsymbol\mu(\partial_b\mathbf I')\boldsymbol\mu. \tag{4.10l}
\end{equation*}
\textbf{Linearity:} $\mathbf I'=\mathbf I_N-\boldsymbol\sigma\boldsymbol\sigma^{\mathsf T}$ with $\boldsymbol\sigma_b=\sum_a\boldsymbol\zeta_{ab}Q_a$, so $\partial_b\mathbf I'$ is built from \emph{constant} $\boldsymbol\zeta$'s. Substituting, the gradient terms cancel and the Coriolis sum rules collapse everything to a single trace:
\begin{equation*}
\boxed{\;\hat U(Q)=-\frac{\hbar^2}{8}\sum_{\alpha=1}^{3}\mu_{\alpha\alpha}(Q).\;}\tag{4.11}
\end{equation*}
\centering\footnotesize A pure $c$-number from the \emph{measure} --- the electrons do not modify it.
\end{frame}

% ==================================================================
\setstage{4}
\section{The complete quantum Hamiltonian}

\begin{frame}{Electronic kinetic operator}
\setstage{4}
In BF Cartesian electron coordinates $\hat{\bar{\mathbf p}}_i=-i\hbar\nabla_{\bar{\mathbf r}_i}$ is already Hermitian on the flat measure, so the electron + mass-polarization piece quantizes directly:
\begin{equation*}
\hat T_e=\sum_i\frac{\hat{\bar{\mathbf p}}_i^{\,2}}{2m_e}+\frac{1}{2M_N}\Bigl(\sum_i\hat{\bar{\mathbf p}}_i\Bigr)^{2}
=\sum_i\frac{\hat{\bar{\mathbf p}}_i^{\,2}}{2\mu_e}+\frac{1}{M_N}\sum_{i<l}\hat{\bar{\mathbf p}}_i\!\cdot\!\hat{\bar{\mathbf p}}_l. \tag{4.12, 4.13}
\end{equation*}
The diagonal part renormalizes the electron mass, $\tfrac1{m_e}+\tfrac1{M_N}=\tfrac1{\mu_e}$; the off-diagonal two-electron term is the mass polarization \emph{sensu stricto} (nuclear recoil).
\end{frame}

\begin{frame}{The exact non-relativistic molecular Hamiltonian}
\setstage{4}
Collecting the rovibrational kernel (3.22), the Watson pseudopotential (4.11), and $\hat T_e$:
\begin{equation*}
\boxed{
\begin{aligned}
\hat H_{\mathrm{rve}}=\;&\tfrac12\!\sum_{\alpha\beta}\bigl(\hat J_\alpha-\hat L_{\mathrm{elec},\alpha}-\hat\pi_\alpha\bigr)\mu_{\alpha\beta}(Q)\bigl(\hat J_\beta-\hat L_{\mathrm{elec},\beta}-\hat\pi_\beta\bigr) &&\text{\footnotesize rotation--vibration--electron}\\[-1pt]
&+\tfrac12\!\sum_b\hat P_b^2 &&\text{\footnotesize vibration}\\[-1pt]
&+\sum_i\frac{\hat{\bar{\mathbf p}}_i^{\,2}}{2m_e}+\frac{1}{2M_N}\Bigl(\sum_i\hat{\bar{\mathbf p}}_i\Bigr)^{2} &&\text{\footnotesize electrons + mass pol.}\\[-1pt]
&-\frac{\hbar^2}{8}\sum_\alpha\mu_{\alpha\alpha}(Q) &&\text{\footnotesize Watson pseudopotential}\\[-1pt]
&+V_{NN}+V_{Ne}+V_{ee} &&\text{\footnotesize Coulomb}
\end{aligned}}\tag{4.14}
\end{equation*}
\centering\footnotesize Hermitian on $L^2(\mathcal C,d\mu_W)$, $\;d\mu_W=\sin\theta\,d\phi d\theta d\chi\prod_b dQ_b\prod_i d^3\bar{\mathbf r}_i$.
\end{frame}

% ==================================================================
\section{Summary}

\begin{frame}{The journey, in four stages}
\setstage{4}
\begin{enumerate}
  \item[\textbf{1.}] \alert{Classical}: lab $H$ (1.1) $\to$ frame-independent $\tilde H$ (2.22), COM split off exactly for any $\mathbf P_{\mathrm{cm}}$.
  \item[\textbf{2.}] \alert{Body-Fixed}: Lagrangian (electrons stay canonical) $\to$ key identity $\mathbf J-\mathbf L_{\mathrm{elec}}=\mathbf I_N\boldsymbol\omega+\mathbf L_N^{\mathrm{vib}}$ $\to$ classical $H_{\mathrm{rve}}$ (3.22) with Watson $\boldsymbol\mu=(\mathbf I_N-\boldsymbol\sigma\boldsymbol\sigma^{\mathsf T})^{-1}$.
  \item[\textbf{3.}] \alert{Quantize}: Podolsky operator (4.3) + rescaling (4.5); anomalous algebra $[\hat J_\alpha,\hat J_\beta]=-i\hbar\epsilon_{\alpha\beta\gamma}\hat J_\gamma$; Watson pseudopotential $\hat U=-\tfrac{\hbar^2}{8}\sum_\alpha\mu_{\alpha\alpha}$.
  \item[\textbf{4.}] \alert{Quantum $\hat H$}: the exact non-relativistic $\hat H_{\mathrm{rve}}$ (4.14).
\end{enumerate}
\end{frame}

\begin{frame}{What we did --- and did not --- assume}
\setstage{4}
$\hat H_{\mathrm{rve}}$ (4.14) is \alert{exact}. The only inputs were:
\begin{itemize}
  \item non-relativistic dynamics;
  \item the Eckart frame (a convention);
  \item $3N_{\mathrm{nucl}}-6$ normal coordinates (an exact parametrization).
\end{itemize}
\vspace{0.5ex}
No Born--Oppenheimer, no clamped nuclei, no neglect of mass polarization, no harmonic assumption.
\begin{block}{Outlook}
(4.14) is the rigorous starting point: a Born--Huang expansion in the parametric electronic basis, then the Born--Oppenheimer truncation, recovers the familiar single-surface picture --- a controlled approximation \emph{of an exact operator}.
\end{block}
\end{frame}

\begin{frame}[standout]
The full molecular Hamiltonian: \\ exact, quantized, and ready to approximate.
\end{frame}

\end{document}
